\section{Evaluation Protocol}
\label{sec:evaluation}
We evaluate translations with a layered protocol that moves from strict syntax, to meaning, to the reliability of the evaluation itself, and finally to deployment efficiency. Every metric below is reported as a mean over independent runs rather than a single measurement.

\subsection{Multi-Seed Aggregation}
LLM decoding and fine-tuning are sensitive to initialization, so each fine-tuned configuration is run over several training seeds on the canonical split, and accuracy is reported as the seed mean together with a $95\%$ confidence interval. This makes the headline numbers stable estimates of capability rather than artifacts of a single favorable run.

\subsection{Exact Match and Meaning}
\noindent\textbf{Exact match.} Predictions are first compared to the reference under a normalization that removes purely surface differences: whitespace and letter case are normalized, logical operators are mapped to a canonical ASCII form (e.g.\ $\land,\lor,\neg,\to$), and the spelling of coalition names inside $\langle\!\langle\cdot\rangle\!\rangle$ is normalized while propositions are left intact, so that a non-semantic agent renaming is not penalized but a genuine lexical change still counts as a miss. Because a single requirement may admit several jointly required readings (for instance under quantifier-scope ambiguity), exact match demands that \emph{all} gold formulas be produced, not merely one of them. This deterministic check is the syntactic-precision baseline.

\noindent\textbf{LLM-as-a-judge.} A prediction that is \emph{not} an exact match is forwarded to an LLM-as-a-judge, which decides strategic and temporal equivalence to the reference and returns a structured verdict (a \texttt{correct} label and a free-text rationale). To avoid depending on any single model, we benchmark \emph{six} independent LLM judges (DeepSeek-V3.2, GPT-4.1, GPT-5.2, GPT-5.4, and the open-weight Gemma-2-27B and Llama-3.3-70B) under one fixed prompt (version~v1.4), whose adjudication rubric and calibration examples are reproduced in the technical appendix, and validate each against the expert audit below. Headline accuracy is reported under the single most human-aligned judge (the open-weight Llama-3.3-70B), which is moreover \emph{not} among the evaluated systems and so never grades its own outputs. Judge decisions are cached per judge identity, so the same judge is never queried twice on an identical triple, while the judges always remain independent. Reported accuracy is therefore the exact-match rate plus the judge-recovered fraction, a decomposition we state explicitly.

\noindent\textbf{Why an LLM judge rather than formal equivalence.} One might instead decide semantic equivalence of two ATL/ATL$^\ast$ formulas over concurrent game structures, but this is ill-posed here: the atomic propositions and agent names are grounded in the natural-language request rather than in a fixed game structure, so there is no shared model over which to evaluate satisfaction, and a faithful prediction may use a differently named but corresponding proposition or coalition; equivalence for ATL$^\ast$ is also computationally heavy. The judge therefore targets the appropriate notion: whether the prediction preserves the \emph{strategic-temporal intent} of the reference (coalition attribution, the strategic modality and its polarity, the temporal operators, and the scope relations among them) rather than verbatim form, requiring that all admissible readings appear for requirements with multiple readings. We treat these verdicts as estimates and validate them against human judgment next.

\subsection{Reliability of the Evaluation}
\noindent\textbf{Inter-rater agreement.} To quantify how trustworthy the automatic verdicts are, we compute pairwise Cohen's $\kappa$, Fleiss' $\kappa$, and Krippendorff's $\alpha$ across judges. Agreement is measured \emph{only} over items that actually reach the judges: deterministic verdicts---exact matches, and empty or unmatched predictions scored as incorrect without an LLM call---are excluded, since including them would trivially inflate agreement.

\noindent\textbf{Human validation.} An expert audit provides ground truth for the judges themselves. A stratified sample of $600$ items---deliberately over-representing items on which the judges disagree---is labeled independently by two expert annotators (inter-annotator agreement $99.5\%$, Cohen's $\kappa=0.99$). The annotators then deliberated over their three initial disagreements and reached agreement on two of them; the third resisted consensus and is retained as a genuine unresolved disagreement, excluded from the human-as-reference comparison and leaving $599$ consensus labels. On both reconciled cases the annotators concluded that the prediction \emph{was} faithful, overturning the stricter initial reading that the automated judges had echoed, so here the human audit corrects an over-strict tendency shared by the LLM judges rather than merely confirming them. The single remaining case sits on the boundary of strategic--temporal equivalence, an inherent limit that no equivalence-labeling protocol, human or automated, can fully avoid. The consensus labels both measure how closely each judge tracks human judgment and produce a robustness ranking in which human labels \emph{override} the judges wherever an item was audited. Because the judges track human judgment unevenly, we report headline accuracy under the most human-aligned judge alone, with two cross-checks for robustness: a conservative mean over DeepSeek-V3.2 and the stricter GPT-5.2 (both external to the evaluated systems), and a mean over all six judges.

\subsection{Efficiency}
We characterize deployment efficiency through the accuracy--latency trade-off rather than a single scalar: for every system we report mean per-example latency alongside accuracy and identify the Pareto frontier, the systems for which no alternative is simultaneously faster and more accurate. We avoid monetary cost models, since token prices and amortized GPU-hour rates are provider- and hardware-specific, whereas accuracy and latency are measured directly. One caveat applies to latency: cloud models are timed end-to-end over the network, so their latency is comparable within the local group but only indicative across the API/local boundary, while accuracy remains comparable throughout.

\subsection{Significance and Transparency}
To establish whether the leading system is genuinely ahead of its closest competitors, we compare systems on shared items with a paired bootstrap confidence interval and a randomization test, rather than the marginal accuracy gap alone. We also disclose every provider content-filter event: the few inputs rejected by the cloud safety filter (at generation or judging time) count as incorrect in the reported accuracy, are listed explicitly, and the ranking is recomputed with them excluded to confirm that the conclusions do not depend on them.



%\paragraph{Architectural Setup.}
%All experiments were conducted using the \texttt{nl2atl} framework implemented in Python~3.10 and executed on a Linux-based environment. The inference backend is built on a modular architecture supporting both local open-weight models and cloud-based APIs, with local inference executed on GPU-enabled nodes equipped with NVIDIA accelerators. Model orchestration, data handling, and evaluation pipelines are implemented using standard scientific Python libraries, while large-scale experimental runs are managed via SLURM on a high-performance computing cluster. The model-checking and syntactic validation of generated ATL/ATL$^\ast$ formulas are performed using the \textsc{VITAMIN} verification tool, ensuring that all evaluated outputs are compatible with an existing strategic model-checking pipeline.



%\subsection{Use Case Scenario} 
%Bob is an engineer working with a CGS in which he must guarantee a winning outcome. To reason about such strategic objectives, he requires a formalism for strategic reasoning with semantics that enforce victory, such as ATL. However, Bob is not an expert in formal logics and lacks the knowledge required to construct well-formed ATL formulas. Therefore, he seeks assistance from an LLM to translate his informal specifications into a formally correct verifiable goal. "In every interaction, the user can, if he is in coalition with the ticket machine, obtain a ticket after paying for it."


%\paragraph{Discussion.}
%This case study illustrates how the specifications produced by our NL-to-ATL$^*$ translator can be grounded and evaluated in simulation-based environments using a finite-trace interpretation of a strategic logic fragment. While this evaluation does not constitute ATL$^*$ model checking in the classical sense, it provides a principled and reproducible methodology for assessing strategic requirements in settings where full game-theoretic models are unavailable. We view this as a complementary validation approach that extends the practical applicability of NL-driven formal specification frameworks to realistic multi-agent simulations.




